\documentclass[11pt,reqno]{amsart}
\usepackage{amsmath,amssymb,graphicx,amscd}
\usepackage{mathrsfs}
\usepackage{enumerate}
% pour les S des spheres
\usepackage{mathrsfs}
\usepackage{dsfont}
\usepackage{stmaryrd}
\usepackage{wrapfig}
\usepackage{pstricks}

% ----------------------------------------------------------------
\vfuzz2pt % Don't report over-full v-boxes if over-edge is small
\hfuzz2pt % Don't report over-full h-boxes if over-edge is small

\newcommand{\prob}{{\mathbb{P}}}

\renewcommand{\S}{\mathscr{S}}  % complex unit sphere

\newcommand{\CC}{{\mathbb C}}
\newcommand{\HH}{{\mathbb H}}
\newcommand{\RR}{{\mathbb R}}
\newcommand{\QQ}{{\mathbb{Q}}}
\newcommand{\NN}{{\mathbb{N}}}
\newcommand{\ZZ}{{\mathbb{Z}}}
\renewcommand{\P}{{\mathbb{P}}}
\renewcommand{\span}{{\mathrm{span}}}
\newcommand{\sgn}{\operatorname{{sgn}}}
 \DeclareMathOperator{\Prob}{\mathbb{P}}   %probability
 \DeclareMathOperator{\E}{\mathbb{E}}      %expectation
 \DeclareMathOperator{\cov}{cov}             %logarithmic integral

 \newcommand{\I}{1\!\!1}                   %Indicator function


 \newcommand{\dd}{{\mathrm{d}}}            %differential
 \newcommand{\ii}{{\mathrm{i}}}
  \newcommand{\Id}{{\mathrm{Id}}}
\newcommand{\law}{\overset{\mbox{\rm \scriptsize law}}{=}}
\newcommand{\convlaw}{\overset{\mbox{\rm \scriptsize law}}{\longrightarrow}}


\renewcommand{\Re}{{\mathfrak{Re}}}
\renewcommand{\Im}{{\mathfrak{Im}}}

% THEOREMS -------------------------------------------------------

\newtheorem{thm}{Theorem}[section]
\newtheorem{cor}[thm]{Corollary}
\newtheorem{lem}[thm]{Lemma}
\newtheorem{prop}[thm]{Proposition}
\theoremstyle{definition}
\newtheorem{defn}[thm]{Definition}
\theoremstyle{remark}
\newtheorem*{rem}{Remark}
 \newcommand{\rd}{{\mathrm{d}}}            %differential
\newcommand{\1}{\mathds{1}}


\numberwithin{equation}{section}



\begin{document}

~\vspace{-3cm}
\begin{center}
{\bf Probability, homework 4 due October 7.}
\end{center}
\vspace{1cm}
\small




\noindent{\bf Exercise 1}.
Let $n$ and $m$ be random numbers chosen independently and uniformly on $\llbracket 1, N\rrbracket$.
What are $\Omega, \mathcal{A}$ and $\P$ (which all implicitly depend on $N$) ?
Prove that
$
\P(n\wedge m=1)\underset{N\to\infty}{\longrightarrow}\zeta(2)^{-1}
$
where $\zeta(2)=\prod_{p\in\mathcal{P}}(1-p^{-2})^{-1}=\sum_{n\geq 1}n^{-2}=\frac{\pi^2}{6}$ (you don't have to prove these equalities). Here $\mathcal{P}$ is the set of prime numbers and $n\wedge m=1$ means that their greatest common divisor is 1.\\


\noindent{\bf Exercise 2}. Let $X$ be a random variable with density $f_X(x)=(1-|x|)\mathds{1}_{(-1,1)}(x)$. Show that its characteristic function is
$$
\phi_X(u)=\frac{2(1-\cos u)}{u^2}.
$$





\noindent{\bf Exercise 3.} 
\begin{enumerate}
\item Prove that $\hat \mu$ is real-valued if and only if $\mu$ is symmetric, i.e. $\mu(A)=\mu(-A)$ for any Borel set $A$\\
\item If $X$ and $Y$ are i.i.d., prove that $X-Y$ has a symmetric distribution.
\end{enumerate}



\noindent{\bf Exercise 4}.
Let $X_\lambda$ be a real random variable, with Poisson distribution with parameter $\lambda$. Calculate the characteristic function of $X_\lambda$. Conclude that $(X_\lambda-\lambda)/\sqrt{\lambda}$ converges in distribution to a standard Gaussian, as $\lambda\to\infty$.\\



\noindent{\bf Exercise 5.}
Assume that the sequence of random variables $(X_n)_{n\geq 1}$ satisfies
$
\E X_n\to 1
$
and $\E X_n^2\to 1$.
Prove that $(X_n)_{n\geq 1}$ converges in distribution. What is the limit?\\



\noindent{\bf Exercise 6}.
Let $(X_n)_{n\geq 1}$, $(Y_n)_{n\geq 1}$ be real random variables, with $X_n$ and $Y_n$ independent for any $n\geq 1$, and assume that $X_n$ converges in distribution to $X$ and $Y_n$ to $Y$, with $X$ and $Y$ independent defined on the same probability space. Prove that $X_n+Y_n$
converges in distribution to $X+Y$.\\





\noindent{\bf Exercise 7}.
Let $X,Y$ be independent and assume that for some constant $\alpha$ we have $\mathbb{P}(X+Y=\alpha)=1$. Prove that $X$ and $Y$ are both constant random variables.\\




\noindent{\bf Exercise 8.}
Let $f,g:\mathbb{R}\to\mathbb{R}$ be nondecreasing measurable functions. Let $\mu$ be a probability measure on $\mathbb{R}$ and assume $f,g,fg\in{\rm L}^1(\mu)$. prove that
$$
\int fg\rd\mu\geq \int f\rd\mu \cdot \int g\rd\mu.
$$



\noindent{\bf Exercise 9.}
Let $(X_n)_{n\geq 1}$ be a sequence of i.i.d. random variables with standard Cauchy distribution  and let
$M_n=\max(X_1,\dots,X_n).$ Prove that $(n M_n^{-1})_{n\geq 1}$ converges in distribution and identify the limit.\\




\noindent{\bf Exercise 10}. Let $(X_i)_{i\geq 1}$ be a sequence of independent random variables, with $X_i$ uniform on $[-i,i]$.  Let $S_n=X_1+\dots+X_n$.
Prove that $S_n/n^{3/2}$ converges in distribution and describe the limit.\\





\noindent{\bf Exercise 11}.  Find a probability distribution $\mu$ of a $\mathbb{Z}$-valued random variable which is symmetric ($\mu(\{i\})=\mu(\{-i\})$ for any $i\in\mathbb{Z}$),  not integrable,  but such that its characteristic function is differentiable at $0$.\\ 


\noindent{\bf Exercise 12.}
Let $X$, $Y$ be i.i.d., with characteristic functions denoted $\varphi_X,\varphi_Y$, and suppose $\E(X)=0$, $\E(X^2)=1$. Assume also that $X+Y$ and $X-Y$ are independent.  
\begin{enumerate}
\item Prove that 
$$
\varphi_X(2u)=(\varphi_X(u))^3\varphi_X(-u)
$$
\item Prove that $X$ is a standard Gaussian random variable.
\end{enumerate}



\end{document} 